% BORRADOR (13.11.2026) de la seccion del caso Lima para el Cap. V. Va despues de "Resultado central: el segundo
% estudio" y antes de "Que explica el resultado". Los PENDIENTE se llenan con lima/salidas/ y lima/alta/salidas/.
%----------------------------------------------------------------------
\section{Caso Lima: el nudo de Reducto con 28 de Julio}
\label{sec:caso-lima}
\suppressfloats[t]

Esta sección lleva el método a un cruce real de Lima. PENDIENTE: respuesta
en dos o tres frases (frente al plan actual, frente al fijo ajustado al
segundo y frente al actuado ajustado, con los dos niveles de demanda).

\noindent\textbf{Qué se midió.} El caso es el nudo donde se cruzan la
Av.~Reducto, la Av.~28 de Julio y la Av.~Vasco Núñez de Balboa, en
Miraflores, junto al puente sobre la Vía Expresa. Todas las luces del
nudo siguen un mismo programa: es un solo cruce, con un solo
controlador. El Anexo~D cuenta cómo se armó el escenario:
\begin{itemize}
  \item \textbf{Red.} La forma del nudo sale de OpenStreetMap; los
    carriles y los giros, de la observación en el sitio.
  \item \textbf{Semáforo.} El plan que opera hoy se midió en el sitio,
    en la hora punta de la tarde: un ciclo de 132~s, con tres estados
    con verde y 34~s con todo en rojo. Todos los controles conservan
    esos 34~s.
  \item \textbf{Demanda.} No hay conteo por movimiento. Se usa el total
    publicado, 4271~veh/h \citep{Miraflores2017_PMU}, repartido en
    proporción al verde y a los carriles que el plan actual da a cada
    movimiento: es el supuesto más favorable para ese plan. Como el
    cruce está calificado como saturado, se añade una demanda alta, un
    20\,\% mayor (5125~veh/h).
  \item \textbf{Controles.} El plan actual, el fijo ajustado al segundo,
    el actuado ajustado y el modelo completo. Todos usan los mismos
    estados, en el mismo orden; solo cambia cuánto dura cada verde.
  \item \textbf{Modelo.} El modelo completo del segundo estudio, sin
    cambios, entrenado tres veces con cada demanda. Aquí es un solo
    agente, que decide cuánto dura cada estado del programa: no hay
    cruces vecinos ni mensajes.
  \item \textbf{Evaluación.} Como en el segundo estudio: llegadas
    aleatorias, dos tandas de 100 semillas, la misma regla de veredicto
    y reglas escritas antes de entrenar.
\end{itemize}
% fuente: anexo_d_osm.tex; lima/REGISTRO.md (reglas del 13.11.2026, 19:30, y segundo nivel de demanda, 19:40);
% lima/programa.py (PLAN, ROJO_TOTAL), lima/demanda.py (TOTAL_PUBLICADO, GRUPOS, DESTINOS; alta = factor 1.2),
% lima/entrenar.sh y lima/entrenar_nivel.sh (opciones del modelo completo, bases 42, 2042 y 4042)

% fuente: secciones/tablas/lima_resultado.tex (lima/tablas_lima.py): lima/salidas/ y lima/alta/salidas/,
% eval_t1_*.csv y eval_t2_*.csv (retraso_total_prom)
\begin{table}[!tbp]
\setstretch{1}
\centering
\footnotesize
\setlength{\tabcolsep}{3pt}
\input{secciones/tablas/lima_resultado}
\caption[Caso Lima: retraso total con cada control]{Caso Lima: retraso total por vehículo con cada control y cambio del modelo completo frente a cada uno, con la demanda publicada y con la demanda alta. Cada cifra se da para las dos tandas de evaluación, de 100 semillas cada una; un cambio negativo es menos retraso con el modelo. El IC95 de cada cambio está en la Tabla~\ref{tab:lima-completa}.}
\label{tab:lima-resultado}
\end{table}

\noindent\textbf{Lectura.} PENDIENTE (tres puntos: frente al plan actual,
frente al fijo ajustado al segundo y veredicto, frente al actuado
ajustado).

% fuente: images/fig_lima_retraso.pdf (lima/tablas_lima.py, figura_retraso): los retrasos de la tanda 1 de la
% Tabla tab:lima-resultado
\begin{figure}[!tbp]
\centering
\includegraphics[width=0.92\textwidth]{fig_lima_retraso.pdf}
\caption[Caso Lima: retraso total por control]{Caso Lima: retraso total por vehículo (s/veh) con el plan actual, el fijo ajustado al segundo, el actuado ajustado y el modelo completo, en la tanda 1 (100 semillas). Un panel por nivel de demanda, cada uno con su escala.}
\label{fig:lima-retraso}
\end{figure}

\noindent\textbf{Cómo aprendió el modelo.} PENDIENTE (una o dos frases con
la Figura~\ref{fig:lima-aprendizaje} y la Tabla~\ref{tab:lima-entrenamientos}:
en cuántos episodios alcanza a las referencias, si los tres
entrenamientos coinciden y qué verdes usa).

% fuente: images/fig_lima_aprendizaje.pdf (lima/tablas_lima.py, figura_aprendizaje): columna val_total de
% resultados_lima_ippo_vecinos_completo_s<base>.csv y salidas/val_referencias.json de cada nivel
\begin{figure}[!tbp]
\centering
\includegraphics[width=\textwidth]{fig_lima_aprendizaje.pdf}
\caption[Caso Lima: aprendizaje del modelo]{Caso Lima: retraso total del modelo en las ocho semillas de validación a lo largo de sus tres entrenamientos, con cada nivel de demanda. Las líneas grises son el plan actual y el fijo ajustado al segundo en esas mismas semillas. La escala vertical corta los primeros episodios, que quedan muy por encima.}
\label{fig:lima-aprendizaje}
\end{figure}

\noindent\textbf{Qué no dice este caso.}
\begin{itemize}
  \item \textbf{No prueba los mensajes.} Con un solo cruce no hay
    coordinación entre vecinos, que es la otra mitad de la pregunta de
    la tesis. Los vecinos de este nudo son los dos semáforos del puente,
    que no están en la red.
  \item \textbf{La demanda es un supuesto.} El total es el publicado,
    pero el reparto entre movimientos no se contó.
  \item \textbf{El retraso simulado no es el del nudo.} La simulación
    no reproduce el bloqueo desde el puente, los cruces dentro del nudo
    ni los peatones (Anexo~D). Sirve para comparar controles entre sí.
\end{itemize}
% fuente: lima/REGISTRO.md, "Limites conocidos de antemano"; anexo_d_osm.tex, "Lo que el escenario no reproduce"
